\documentclass[11pt]{article}

\usepackage[margin=1in]{geometry}
\usepackage{amsmath}
\usepackage{siunitx}
\usepackage{xcolor}
\usepackage{hyperref}

\hypersetup{
  colorlinks=true,
  linkcolor=blue,
  urlcolor=blue
}

\newcommand{\code}[1]{\texttt{#1}}
\newcommand{\file}[1]{\texttt{#1}}

\title{Structure of \file{full\_analysis\_example.py}}
\author{CASPErG LF 2022-12-23 Milky-Way Axion Analysis}
\date{\today}

\begin{document}
\maketitle

\section{Purpose}

This note describes the analysis stages in
\file{examples/LF-science-runs/MW-axion-analysis-20221223/full\_analysis\_example.py}.
The script is a readable analysis skeleton for the LF 2022-12-23 data set. It
is organized around the pipeline described in \file{\_internal/PRD.md}: device
HDF5 files are loaded by a device-specific \code{Signal} subclass, converted to
the common \code{TS\_raw}/\code{TS} data model, characterized with one-pulse
and NoPulse analysis, searched for axionlike excess power, and finally combined
into candidate and sensitivity/exclusion records.

The implementation details referenced here are primarily in
\file{CASPErG/DarkMatterAnalysis.py}, especially the \code{dmScanStep} and
\code{dmScan} methods used by the script.

\subsection{0. Message Prefix}

The first runtime block defines a single message prefix,
\code{msgPrefix = "[20221223 MW axion analysis]"}. This is used for printed
progress messages so that output from this script can be distinguished from
messages emitted by library methods.

\subsection{1. User Configuration}

This section collects all user-facing choices that define the run: the data
root, output directory, Zurich Instruments device and demodulator identifiers,
the three scan folders, the AxionBloch projection convention, and the numerical
analysis settings.

The sample is defined as C-12 methanol using AxionBloch's \code{Sample} class.
The script also defines values controlling NoPulse segmentation, drift checks,
baseline handling, local-noise estimation, optimal-filter stepping, outlier
thresholding, and candidate consistency. In the PRD language, these parameters
belong to the science-run configuration rather than to the raw HDF5 files.

\subsection{2. Small Helpers for File Discovery}

The helper functions convert the directory structure into structured file
metadata. \code{natural\_key} sorts filenames numerically, so
\code{AxionScan\_2.h5} precedes \code{AxionScan\_10.h5}. \code{trailing\_file\_index}
extracts the scan-step index from a filename. \code{discover\_scan\_files}
collects AxionScan, OnePulse, and CPMG files in one scan folder, and
\code{group\_scan\_files\_by\_index} groups them by their trailing index.

\code{build\_step\_info} converts each file group into dictionaries accepted by
\code{dmScan.generateStepPoolFromFileInfo}. These dictionaries preserve
\code{scanIndex}, \code{scanName}, \code{stepIndex}, and the corresponding
NoPulse, OnePulse, and CPMG file paths. The remaining helper functions are for
runtime reporting and chronological step ordering after NoPulse timestamps are
loaded.

\subsection{3. Build One \code{dmScan} Containing All Three Scans}

This section discovers files from every configured scan folder and combines
them into one \code{dmScan}. This follows the PRD requirement that a science run
with repeated passes over the same band can be represented by one scan object
whose steps carry \code{scanIndex} and \code{stepIndex} metadata.

The call to \code{generateStepPoolFromFileInfo(..., prepare=False)} creates
metadata-only \code{dmScanStep} objects. Each step receives file-backed
\code{ZI\_5MHz\_LIA\_Signal} instances for the available measurement types, but
the full time series is not loaded. This keeps the construction fast and
auditable: paths are attached, normalized file groups are stored, and the
active science-analysis pool remains empty until preparation.

\subsection{4. Prepare Scan Steps and Sort by Measurement Time}

\code{science\_scan.prepareMetadataSteps()} promotes metadata-stage steps into
the active analysis pool only after preparation checks pass. Internally,
\code{dmScan.addStep} calls \code{dmScanStep.prepareForScan}, which checks file
coverage through the attached NoPulse signal without running the scientific
FFT or outlier pipeline.

Steps that fail coverage or readiness checks are retained in
\code{science\_scan.rejectedSteps} together with their preparation reports.
Ready steps enter \code{science\_scan.pool} and are sorted by measurement start
time. This implements the PRD distinction between metadata construction and
scientific analysis.

\subsection{5. Attach Sample, Station, and Axion Configuration}

For each active step, the script attaches the sample, station, and axion
projection metadata needed by later AxionBloch calls. The sample supplies
material and spin information, while \code{Mainz}, \code{axionCase},
\code{biasFieldAxis}, and \code{projectionAxis} define the station and
projection convention used for the axion wind response.

This block intentionally keeps some hardware objects as placeholders. The PRD
requires the pipeline to store such physical metadata explicitly because the
same NoPulse spectrum can have different physical interpretations under
different sample, magnet, station, or projection assumptions.

\subsection{6. Load One-Pulse and NoPulse Data}

This section loads the OnePulse and NoPulse HDF5 files using
\code{ZI\_5MHz\_LIA\_Signal.loadData}. The loader is responsible for translating
the Zurich Instruments HDF5 layout into the common signal contract:
\[
  \mathrm{TS\_raw} = X + iY,
\]
with an organized one-row \code{TS} by default.

After NoPulse loading, the script reconstructs the absolute NoPulse start and
stop times for each file index. These times are copied into both the signal
metadata and the step-level measurement interval. The steps are later sorted
again by this loaded NoPulse timing information.

\subsection{7. OnePulse Stage: Characterize the NMR Line}

The OnePulse stage characterizes the free-decay measurement. The script sets
the one-pulse acquisition delay and duration, then calls
\code{step.characterizeOnePulse(...)}.

In \file{DarkMatterAnalysis.py}, \code{characterizeOnePulse} finds pulses,
Fourier-transforms each organized shot, and fits the mean one-pulse PSD with a
Lorentzian through \code{onePulseSignal.fitPSD\_legacy}. The fit stores the
Larmor frequency on \code{onePulseSignal.LarmorFreq}, stores the NMR linewidth
as \code{step.NMR\_FWHM\_freq}, and derives
\[
  T_2^* = \frac{1}{\pi\,\mathrm{FWHM}_{\mathrm{NMR}}}.
\]
The raw frequency grid and shot PSD remain owned by
\code{step.onePulseSignal.freqs} and \code{step.onePulseSignal.PSD}. The
canonical physical 90-degree response is not stored as a scan-step mirror.
Instead, \code{characterizeOnePulse} extracts \(\Phi_{90\deg}\) from the
integral of the mean one-pulse PSD, corrects for the acquisition-averaged
decay factor,
\[
  \frac{1}{T}\int_{t_0}^{t_0+T}\exp(-2t/T_2^*)\,dt
  =
  \frac{T_2^*}{2T}
  \left[
    \exp(-2t_0/T_2^*)-\exp(-2(t_0+T)/T_2^*)
  \right],
\]
normalizes the measured one-pulse lineshape to unit frequency integral, and
stores
\[
  \code{onePulseSignal.PSD\_90deg}(f)
  =
  \Phi_{90\deg}^2\,\mathcal{L}_{90}(f).
\]
\code{onePulseSignal.PSD\_90deg} is two-dimensional, with the same shape as
\code{onePulseSignal.PSD}, so later step methods can use it directly while the
raw per-shot PSD remains available as a diagnostic.

\subsection{8. Create or Update AxionBloch Axion Models}

After the Larmor fit, the script creates a \code{MilkyWayAxionHalo} object for
each analyzed step using the fitted resonance frequency. This object supplies
the axion linewidth, coherence information, station-dependent lineshape, and
Rabi-frequency response used in later filtering and sensitivity calculations.

The surrounding comments mark this as an interface placeholder. The PRD fixes
the physical convention: use the Milky-Way axion model, use
\code{case="grad\_perp"}, use the zenith bias-field convention, and evaluate the
station-dependent lineshape at the midpoint of the NoPulse measurement rather
than time-averaging over the measurement interval.

\subsection{9. NoPulse Stage: Organize Time Series and Run Sanity Checks}

The NoPulse stage ensures that the search time series is in the common
two-dimensional \code{TS} representation. The script reshapes the loaded raw
series into one segment:
\[
  \code{TS.shape} = (1, N).
\]

It then calls \code{step.runTSSanityChecks}. The implementation checks drift
first, then jumps and Gaussianity on drift-approved rows. The method records
whether the time series was modified, stores the sanity report on both the step
and the signal, and returns a status such as \code{passed},
\code{pausedDrift}, \code{pausedJumpCluster}, or \code{pausedNonGaussian}.
The script currently leaves jump-mask application as a TODO because masking
requires explicit approval and downstream calibration.

After sanity checks, the active pool is sorted by the loaded NoPulse
measurement start time, and step indices are reassigned to match the
chronological pool order.

\subsection{10. Fourier Transform, Baseline Handling, and Search Spectrum}

This section transforms the NoPulse time series into frequency space by calling
\code{step.noPulseSignal.getFS}. The NoPulse PSD and frequency grid are then
used by \code{step.buildAxionTemplate} to build an RBW-binned AxionBloch
template centered on the fitted Larmor frequency.

\code{buildAxionTemplate} uses the NoPulse measurement midpoint, the station,
the axion case, and the requested projection axis to call AxionBloch's
station-dependent lineshape method. It normalizes the template so that its RBW
sum has unit area.

The script then calls \code{step.prepareSearchSpectrum}. This method selects a
region of interest around the fitted Larmor frequency, estimates a baseline
for each PSD row, forms
\[
  R_j(f) = \frac{\mathrm{PSD}_j(f)}{B_j(f)},
  \qquad
  \mathrm{NPE}_j(f) = R_j(f)-1,
\]
and checks baseline flatness before and after correction. The averaged
normalized power excess, \code{NPE\_search}, becomes the input to the power
statistic.

\subsection{11. Axion Template and Optimal Filtering}

The script calls \code{step.optimalFilter(filterStepX=OPTIMAL\_FILTER\_STEP\_X)}
to convert the normalized excess spectrum into an axion-search statistic.
The method supports two regimes. If the axion linewidth is smaller than the
NoPulse RBW, the signal is treated as unresolved and each bin is converted
directly to excess power. If the axion linewidth is resolved, the
RBW-normalized AxionBloch template is used as the matched-filter kernel and is
shifted across the search spectrum in steps of \code{filterStepX} axion
linewidths.

The method stores the filtered frequency grid, excess-power estimate,
global and local noise estimates, filter stride, realized step size, and
filtered normalized power excess \code{filteredNPE}. The example injects a few
large artificial excursions into \code{filteredNPE} for debugging and then
calls \code{plotOptimalFilterResult}, which plots the axion template, the mean
of \code{onePulseSignal.PSD\_90deg}, raw PSD plus baseline, and filtered
statistic.

\subsection{12. Outlier Finding Inside Each Scan Step}

Each step calls \code{findOutliers}. This method searches the filtered scalar
statistic for positive threshold crossings. The configured scalar threshold is
a placeholder until a Monte Carlo threshold module supplies the approved
95-percent signal-pickup threshold.

\code{findOutliers} creates ordinary \code{Outlier} records that retain the
frequency, support region, excess power, uncertainty, baseline information,
RBW, and provenance. If a jump mask has been applied, outlier use is rejected
unless both transfer and statistical calibration records are present.

\subsection{13. Cross-Step and Cross-Scan Candidate Evaluation}

The science-run object then calls \code{science\_scan.evaluateOutliers}. This
method applies the PRD's single outlier-validation rule: an axionlike signal
must appear at a consistent absolute frequency and with a consistent inferred
coupling strength in every step where it should be detectable.

The method groups outliers by frequency tolerance, infers couplings for
frequency-matched members, tests coupling consistency with an inverse-variance
common-value chi-squared calculation, and checks whether sufficiently sensitive
covered steps are missing a required match. Groups that pass become candidate
groups. Failed groups are rejected, and inconclusive groups are retained for
manual review.

\subsection{14. Sensitivity and Exclusion}

The final active analysis block estimates candidate-blind step-level
sensitivity and then builds the scan-level exclusion from that sensitivity.
The script calls \code{science\_scan.calculateAndSaveStepSensitivities(...)},
which creates a new output folder for per-step limit text files. Each step
calls \code{dmScanStep.calculateSensitivityFromNoiseGrid(...)} and writes a
file named like
\code{20221223-MW-axion-analysis+<stepHash>.txt}.

\code{calculateSensitivityFromNoiseGrid} uses the one-pulse grid and the mean
of \code{onePulseSignal.PSD\_90deg} as the NMR response. The raw
\code{onePulseSignal.PSD} remains the source of per-bin shot-to-shot
diagnostic spread. The method interpolates the local filtered-noise estimate
onto the one-pulse response grid, truncates the response far from the NMR
line, obtains the AxionBloch Rabi-frequency response per unit coupling, and
converts a noise threshold into a coupling limit.

The response model follows the PRD:
\[
  \theta_{a,\mathrm{rms}}
  =
  \Omega_{\mathrm{rms}} T_2^*
  \sqrt{\frac{\tau_a}{\tau_a + T_2^*}},
  \qquad
  \mathrm{PSD}_a(f)
  =
  \code{PSD\_90deg}(f)\,\sin^2\theta_{a,\mathrm{rms}}.
\]
Because this response scales as \(g_{aNN}^2\), the script solves backward from
the local threshold power to a limit on \(g_{aNN}\).

The returned sensitivity records are combined by
\code{science\_scan.combineSensitivities(...)} and plotted by
\code{science\_scan.plotSensitivity(...)}. Exclusion is then calculated
separately by \code{science\_scan.calculateExclusion(...)}, which starts from
the combined sensitivity and applies candidate information only in the
exclusion path. \code{science\_scan.plotExclusion(...)} writes the scan-level
exclusion plot and text output.

\subsection{Coverage Plot TODO}

After the sensitivity block, the script carries the PRD TODO for a future
\code{dmScan.plotCoverage(...)} method. The desired diagnostic plot should show
each scan step's sensitive interval centered on its fitted Larmor frequency,
label scan and step identities, and visually down-weight the edges where the
NMR response is smaller.

This diagnostic does not replace candidate consistency or exclusion logic. It
is intended to help analysts inspect the repeated scan coverage over the
searched frequency band.

\subsection{15. Save Reports}

The last block now performs practical size inspection and serialization. It
calls \code{showLargestAttrSizes} on the scan, each analyzed step, and the
associated NoPulse and OnePulse signals. This is intended to explain why a
small top-level \code{dmScan} object can still produce a much larger pickle
when it contains large nested signal arrays and cached spectral products.

The script then calls \code{science\_scan.saveToPkl(...)} to serialize the
current analysis object. The per-step sensitivity text files are already
written in Section 14, while final metadata reports, step summaries, outlier
tables, and candidate tables remain future reporting outputs.

\end{document}
